\documentclass[10pt,a4paper]{amsart}
\usepackage[margin=19mm]{geometry}
\usepackage{amsmath,amssymb,mathtools}
\usepackage[scr=rsfs,cal=euler]{mathalfa}
\usepackage{xfrac}
\usepackage[unicode,hidelinks,bookmarks=false]{hyperref}
\newcommand{\ie}{i.e.\ }
\newcommand{\cf}{cf.\ }
\newcommand{\resp}{respectively\ }
\newcommand{\Subsection}{\subsection}
\providecommand{\nobreakdash}{\nobreak\hbox{-}}
\setlength{\emergencystretch}{2em}
\multlinegap0pt
\usepackage{amssymb}
\usepackage{mathtools}
\usepackage{bm}
\usepackage{xcolor}
\usepackage[shortlabels]{enumitem}
\newtheorem{assumption}{Assumption}
\newtheorem{proposition}{Proposition}
\newcommand{\dd}{\mathrm{d}}
\title{Mean--Covariance Turnpikes in Wasserstein Distributionally Robust Linear-Quadratic Control}
\author{Yanzhi Wu, Zhengping Ji}
\date{Source: arXiv:2608.26986v1 (CC BY 4.0)}
\begin{document}
\maketitle

\noindent\emph{Source and licence.} Mean--Covariance Turnpikes in Wasserstein Distributionally Robust Linear-Quadratic Control. Version arXiv:2608.26986v1 (CC BY 4.0), distributed by arXiv under the Creative Commons Attribution 4.0 International licence, http://creativecommons.org/licenses/by/4.0/, as recorded in the arXiv licence metadata for this exact version and shown on the abstract page https://arxiv.org/abs/2608.26986v1. Full licence notice: \texttt{licenses/arxiv-2608.26986-CC-BY-4.0.txt}.\par\smallskip

\noindent\textit{Editorial scope.} Counted unit: the article's proposition prop:static-equilibrium (source0.tex lines 905--921). The article's complete original proof of it is delivered with it (source0.tex lines 923--987, one \texttt{proof} environment of 65 lines). No mathematics is written, repaired or re-derived: the statement and the proof are the article's own lines, in the article's own notation, and no retained line is substituted or re-flowed.  Retained uncounted as the source-local context the proof needs: lines 199--203; lines 212--234; lines 242--257; lines 491--494; lines 499--502; lines 702--726; lines 854--859; lines 872--877; lines 881--902.  Nothing was omitted from any retained span, so the package carries no omission record.  No retained line contains a citation, so the wrapper carries no bibliography and no bibliography entry.  No retained line contains a figure environment, an image-inclusion command or drawing code, so no image is delivered and the package declares no asset dependency.  Editorial formatting only, in addition: an AMS \texttt{amsart} preamble that restates the article's own declarations and macros used by the retained lines, namely the environments \texttt{assumption}, \texttt{proposition} on separate counters, together with the standard packages those lines need.  The retained environments print in this document as Assumption 1, Proposition 1 in order of appearance, whereas the article prints the same items with the numbering of its own full document.  The complete native source of the article as distributed in the arXiv e-print for this exact version is bundled unchanged as \texttt{sources/208675/source0.tex}.
\begin{equation}
	x_{t+1}=Ax_t+Bu_t+\Xi w_t,
	\qquad t\ge 0,
	\label{eq:sys}
\end{equation}

The probability distribution of \(w_t\) is unknown.  Instead, at each stage \(t\), a finite disturbance dataset
\(
\widehat w_t^{(1)},\ldots,\widehat w_t^{(M)}\in\mathbb R^k
\)
is available.  Here \(\widehat w_t^{(i)}\) denotes the \(i\)-th empirical disturbance sample at stage \(t\).  Let \(\delta_a\) denote the Dirac probability measure concentrated at \(a\).  The corresponding empirical distribution is
\begin{equation}
	\nu_t:=\frac1M\sum_{i=1}^M\delta_{\widehat w_t^{(i)}} .
	\label{eq:empirical-law}
\end{equation}
 Its empirical mean and covariance are
\begin{equation}
	\begin{aligned}
		d_t
		&:=
		\frac1M\sum_{i=1}^M\widehat w_t^{(i)},\\
		\Sigma_t
		&:=
		\frac1M\sum_{i=1}^M
		(\widehat w_t^{(i)}-d_t)
		(\widehat w_t^{(i)}-d_t)^\top .
	\end{aligned}
	\label{eq:mean-cov}
\end{equation}

\begin{assumption}
	\label{ass:stationary-empirical}
	The empirical mean and covariance in \eqref{eq:mean-cov} are independent of \(t\), i.e.,
	 $d_t\equiv d,
		\Sigma_t\equiv\Sigma$. The initial state $x_0$ in \eqref{eq:sys} has distribution \(\mu_0\in\mathcal P_2(\mathbb R^n)\), with mean and
	covariance
	\begin{equation}
		\begin{aligned}
			m_0:=\!\!\int \!x\,\dd\mu_0(x),
			~X_0\!:=\!\!\int \!(x-m_0)(x-m_0)^\top\dd\mu_0(x).
		\end{aligned}
		\label{eq:initial-ensemble-moments}
	\end{equation}
	The deterministic initial state is recovered by taking
	\(\mu_0=\delta_{x_0}\).
\end{assumption}

\begin{equation}
	W:=BR^{-1}B^\top-\lambda^{-1}\Xi\Xi^\top .
	\label{eq:W-def-bellman}
\end{equation}

\begin{assumption}
	\label{ass:psd-W}
	$W\succeq0$, the pair \((A,W^{1/2})\) is stabilizable, and $(A,Q^{1/2})$ is observable.
\end{assumption}

\begin{proposition}
	\label{prop:riccati-convergence}
	Define the Riccati map
\begin{equation}
	\mathcal R(P):=Q+A^\top(I+PW)^{-1}PA .
	\label{eq:W-R-map}
\end{equation}
with $W$ defined in \eqref{eq:W-def-bellman}. Starting from \(Q_f\), define
\begin{equation}
	S_0:=Q_f,\qquad
	S_{j+1}:=\mathcal R(S_j).\qquad j\ge0
	\label{eq:S-orbit}
\end{equation}
Then under Assumption~\ref{ass:psd-W}, the sequence \eqref{eq:S-orbit} is well defined, remains positive semidefinite and bounded, and converges to a stabilizing positive-semidefinite solution
	\(P_{\rm ss}\) of
	\(
	P=\mathcal R(P).
	\)
	Moreover,
	\begin{equation}
		A_c:=(I+WP_{\rm ss})^{-1}A
		\label{eq:Ac-def}
	\end{equation}
	is Schur.
\end{proposition}

\begin{equation}
	u=-R^{-1}B^\top p,
	\qquad
	\bar w=d+\lambda^{-1}\Xi^\top p .
	\label{eq:stationary-u-w}
\end{equation}

\begin{equation}
	\min_{m\in\mathbb R^n}
	\max_{p\in\mathbb R^n}
	\mathscr H_{\rm s}(m,p).
	\label{eq:static-minimax}
\end{equation}

The stationarity conditions of
\eqref{eq:static-minimax} are
$p=Qm+A^\top p$, $m=Am-Wp+\Xi d$.
By \eqref{eq:stationary-u-w}, the second equation in
\eqref{eq:reduced-static-first-order} recovers the static mean relation
stated above. Equivalently,
\begin{equation}
	\underbrace{
		\begin{bmatrix}
			I-A & W\\
			-Q & I-A^\top
		\end{bmatrix}
	}_{\mathcal K_e}
	\begin{bmatrix}
		m\\p
	\end{bmatrix}
	=
	\begin{bmatrix}
		\Xi d\\0
	\end{bmatrix}.
	\label{eq:static-KKT-matrix}
\end{equation}

\begin{proposition}
	\label{prop:static-equilibrium}
	Under Assumptions~\ref{ass:stationary-empirical} and
	\ref{ass:psd-W}, the reduced static Hamiltonian problem
	\eqref{eq:static-minimax} has a unique saddle point
	\((m^e,p^e)\). It is characterized by
	\eqref{eq:static-KKT-matrix}. The associated static control input and
	static mean of the worst-case disturbance distribution are
	\begin{equation}
		u^e=-R^{-1}B^\top p^e,
		\qquad
		\bar w^e=d+\lambda^{-1}\Xi^\top p^e.
		\label{eq:ue-we}
	\end{equation}
	If \(d=0\), then
	$(m^e,p^e,u^e,\bar w^e)=(0,0,0,0)$.
\end{proposition}

\begin{proof}
Since \(Q\succeq0\) and \(W\succeq0\),
	\(\mathscr H_{\rm s}\) is convex in \(m\) and concave in \(p\). Its
	first-order conditions are
	\begin{equation}
		Qm+(A^\top-I)p=0,
		\qquad
		(A-I)m-Wp+\Xi d=0,
		\label{eq:reduced-static-first-order}
	\end{equation}
	which are equivalent to \eqref{eq:static-KKT-matrix}.
	
	Proposition~\ref{prop:riccati-convergence} gives the stabilizing Riccati
	solution \(P_{\rm ss}\) and the Schur matrix
	\(A_c=(I+WP_{\rm ss})^{-1}A\). The Riccati equation and the identities
	\(
	P_{\rm ss}(I+WP_{\rm ss})^{-1}
	=(I+P_{\rm ss}W)^{-1}P_{\rm ss}
	\)
	and
	\(
	I-W(I+P_{\rm ss}W)^{-1}P_{\rm ss}
	=(I+WP_{\rm ss})^{-1}
	\)
	yield
	\begin{equation*}
		\begin{aligned}
			&
			\begin{bmatrix}
				(I+WP_{\rm ss})^{-1}&0\\
				A^\top P_{\rm ss}(I+WP_{\rm ss})^{-1}&I
			\end{bmatrix}
			\mathcal K_e
			\begin{bmatrix}
				I&0\\
				P_{\rm ss}&I
			\end{bmatrix}\\
			&\qquad=
			\begin{bmatrix}
				I-A_c&(I+WP_{\rm ss})^{-1}W\\
				0&I-A_c^\top
			\end{bmatrix}.
		\end{aligned}
	\end{equation*}
	The two matrices multiplying \(\mathcal K_e\) are nonsingular. The matrix on
	the right-hand side is block upper triangular, and its diagonal blocks are
	nonsingular because \(A_c\) is Schur. Hence \(\mathcal K_e\) is nonsingular,
	and \eqref{eq:static-KKT-matrix} has a unique solution.
	
For \(m,p\in\mathbb R^n\), the first-order conditions at
	\((m^e,p^e)\) give
	\begin{align*}
		&\mathscr H_{\rm s}(m,p^e)-\mathscr H_{\rm s}(m^e,p^e)=(m-m^e)^\top Q(m-m^e)\ge0,\\
		&\mathscr H_{\rm s}(m^e,p)-\mathscr H_{\rm s}(m^e,p^e)=-(p-p^e)^\top W(p-p^e)\le0.
		\label{eq:static-saddle-max}
	\end{align*}
	Therefore, $\mathscr H_{\rm s}(m^e,p)
			\le \mathscr H_{\rm s}(m^e,p^e)\le \mathscr H_{\rm s}(m,p^e)$ for all
			$m,p\in\mathbb R^n$.
	Thus \((m^e,p^e)\) is a saddle point. Every saddle point satisfies
	\eqref{eq:reduced-static-first-order}; uniqueness follows from the
	nonsingularity of \(\mathcal K_e\). Equation~\eqref{eq:ue-we} uniquely
	determines \(u^e\) and \(\bar w^e\). If \(d=0\), uniqueness gives the zero
	solution.
\end{proof}

\end{document}
